\documentclass[11pt]{article}
\usepackage[a4paper,margin=1in]{geometry}
\usepackage{fontspec}
\usepackage[boldfont=false]{xeCJK}
\setCJKmainfont{FandolSong-Regular.otf}
\usepackage{amsmath}
\usepackage{amssymb}
\usepackage{array}
\usepackage{booktabs}
\usepackage{graphicx}
\usepackage{adjustbox}
\usepackage{enumitem}
\setlist[itemize]{topsep=2pt,partopsep=0pt,parsep=0pt,itemsep=2pt,leftmargin=1.4em}
\setlist[enumerate]{topsep=2pt,partopsep=0pt,parsep=0pt,itemsep=2pt,leftmargin=1.6em}
\usepackage{parskip}
\usepackage{fvextra}
\fvset{breaklines=true,breakanywhere=true,fontsize=\small}
\setlength{\parindent}{0pt}

\begin{document}

\begin{center}
  {\LARGE\bfseries Reaction kinetics}\\[0.35em]
  {\large A Level Chemistry}
\end{center}
\vspace{0.6em}

\section*{Rate equations and orders}

A \textbf{rate equation} 速率方程 shows how the rate depends on the concentrations of the reactants:

\[
\text{rate} = k\,[\text{A}]^m[\text{B}]^n
\]

\begin{itemize}
\item $m$ is the \textbf{order of reaction} 反应级数 with respect to A, and $n$ the order with respect to B. Each is $0$, $1$ or $2$.

\item the \textbf{overall order of reaction} 总反应级数 is $m + n$.

\item $k$ is the \textbf{rate constant} 速率常数. The rate equation can only be found by experiment, not from the balanced equation.

\end{itemize}

\subsection*{Finding the order}

\begin{itemize}
\item \textbf{initial rates method}: change one concentration at a time and see how the starting rate changes. If doubling $[\text{A}]$ doubles the rate, the order in A is 1; if it quadruples the rate, the order is 2; if the rate is unchanged, the order is 0.

\item \textbf{graphs}: a concentration–time graph for a first-order reaction has a constant \textbf{half-life} 半衰期 (the time for the concentration to halve). A rate–concentration graph is a straight line through the origin for first order, and a curve for second order.

\end{itemize}

\subsection*{Half-life and the rate constant}

For a \textbf{first-order} reaction the half-life is \textbf{constant} — it does not depend on the concentration. You can find the rate constant from it:

\[
k = \frac{0.693}{t_{\frac{1}{2}}}
\]

You can also find $k$ by putting initial-rate data into the rate equation.

\section*{Reaction mechanisms}

Most reactions happen in several steps. The slowest step is the \textbf{rate-determining step} 决速步骤, and it controls the overall rate. Only the species involved up to and including this step appear in the rate equation.

\begin{itemize}
\item an \textbf{intermediate} 中间体 is a species made in one step and then used up in a later step. It is not in the overall equation.

\item you can suggest a mechanism that fits both the rate equation and the overall equation, predict the order from a given mechanism, or pick out the rate-determining step.

\end{itemize}

If you compare the \textbf{initial rate} 初始速率 of different mixtures, you can deduce the rate equation, and from that work out the mechanism.

\subsection*{Effect of temperature}

Raising the temperature increases the rate constant $k$ (more molecules pass the activation energy), so the rate goes up.

\section*{Catalysts}

\textbf{Catalysts} 催化剂 can be homogeneous or heterogeneous.

\subsection*{Heterogeneous catalysts}

A \textbf{heterogeneous catalyst} 多相催化剂 is in a different physical state from the reactants (usually a solid with gases). It works in three stages:

\begin{enumerate}
\item \textbf{adsorption} 吸附: reactant molecules stick to the catalyst surface.

\item the bonds in the reactants are weakened, so they react more easily.

\item \textbf{desorption} 脱附: the product molecules leave the surface.

\end{enumerate}

Examples are iron in the Haber process, and platinum, palladium and rhodium in a catalytic converter.

\subsection*{Homogeneous catalysts}

A \textbf{homogeneous catalyst} 均相催化剂 is in the same physical state as the reactants. It is used up in one step and then reformed in a later step, so it comes back unchanged. Examples are oxides of nitrogen helping to oxidise atmospheric sulfur dioxide, and $\text{Fe}^{2+}$ or $\text{Fe}^{3+}$ speeding up the reaction between $\text{I}^-$ and $\text{S}_2\text{O}_8^{2-}$.


\end{document}
